\documentclass{article}
\usepackage{mathpartir}

\begin{document}

\begin{inference-notation}
infix:53 " ∈ " => "\in"
infix:53 " ∉ " => "\notin"
infix:53 " ≠ " => "\ne"
infix:53 " = " => "="
infixr:35 " ∧ " => "\land"
infix:50 " → " => "\mathrel{\to_{\beta}}"
infix:50 " ↠ " => "\mathrel{\to_{\beta}^{*}}"
notation:56 "λ" x:atom "." e:56 => "\lambda " x ".\," e
notation:75 "FV" e:76 => "\mathrm{FV}\," e
notation:53 z:atom "fresh" => z "\;\mathsf{fresh}"
notation:85 e:85 "[" s:0 "/" x:atom "]" => e "[" s "/" x "]"
\end{inference-notation}

Variables range over names, and terms are identified up to alpha-equivalence.

\begin{inference}[kind=bnf,rule_set=Untyped lambda syntax]
--------------- Terms
e,f,s ::= x | e₁ e₂
        | λ x . e
\end{inference}

Substitution replaces free occurrences and avoids capture. In the guards,
\(\mathrm{FV}\,s\) denotes the free variables of \(s\). In the final equation,
\(z\;\mathsf{fresh}\) means that \(z\) occurs nowhere in \(e\) or \(s\) and
is distinct from \(x\) and \(y\).

\begin{inference}[kind=equation,rule_set=Capture-avoiding substitution]
x [s/y] = s -- if x = y
x [s/y] = x -- if x ≠ y

(e₁ e₂) [s/x] = (e₁ [s/x]) (e₂ [s/x])

(λ x . e) [s/x] = λ x . e
(λ y . e) [s/x] = λ y . e [s/x] -- if y ≠ x ∧ y ∉ FV s
(λ y . e) [s/x] = λ z . e [z/y] [s/x] -- if y ≠ x ∧ y ∈ FV s ∧ z fresh
\end{inference}

\begin{inference}[rule_set=Beta reduction]
──────────────────────────── R-Beta
(λ x . e) s → e [s/x]

e → f
──────────────────────────── R-AppLeft
e s → f s

e → f
──────────────────────────── R-AppRight
s e → s f

e → f
──────────────────────────── R-UnderBinder
λ x . e → λ x . f
\end{inference}

\begin{inference}[rule_set=Multi-step reduction]
──────────────────────────── R-ManyRefl
e ↠ e

e → f
f ↠ g
──────────────────────────── R-ManyStep
e ↠ g
\end{inference}

\end{document}
